\section*{Contribution and Novelty}
\label{sec:contribution-novelty}
\paragraph{Contribution and Novelty.}
This work establishes a smoothly regularized three-field continuum framework for bone remodeling that integrates apparent density, logarithmic fabric anisotropy, and a nonlocal mechanotransductive stimulus into a single coupled system. A central methodological innovation is the numerically robust treatment of fabric-guided orthotropy via stabilized Sylvester spectral projectors. Unlike standard eigensolver-based approaches, this formulation smoothly blends the spectral decomposition to the isotropic and transversely isotropic limits, thereby enforcing partition of unity and preventing basis flips or discontinuous updates near eigenvalue degeneracy. This ensures that the constitutive law and fabric evolution remain well-defined and differentiable throughout the simulation, even in the presence of evolving anisotropy.

From a computational perspective, we address the convergence challenges inherent to stiff remodeling kinetics by introducing a partitioned block Gauss-Seidel solver stabilized by Anderson (Pulay/DIIS) acceleration. We demonstrate that this acceleration scheme, augmented with adaptive restarts and step limiting, delivers robust convergence in regimes where classical Picard iteration fails. Furthermore, the model introduces a compartment-aware reference stimulus that adapts to local tissue types, which we show to be essential for reproducing physiological density distributions in a population-averaged proximal femur study. Together, these contributions provide a stable, mesh-independent, and extensible platform for predictive bone mechanobiology.
